\documentclass{ximera}
\title{Vector bundles and trivializations}

\begin{document}

    \begin{abstract} Vector bundles, trivializations, basic constructions, basic properties. \end{abstract}
    \maketitle

A reference for this material is Chapter 10  of John M. Lee. Introduction to smooth manifolds. Second edition. Grad. Texts in Math., 218. Springer, New York, 2013. xvi+708pp. ISBN: 978-1-4419-9981-8. 

\begin{definition}
 Let $M$ be a smooth manifold and $k \in \mathbb{N}$. The \emph{trivial vector bundle}  of rank $k$ over $M$ is the product manifold $M \times \mathbb{R}^k$ equipped with the projection $\pi : M \times \mathbb{R}^k \to M $. 

 A \emph{section} of this vector bundle is a smooth map $s : M \to  M \times \mathbb{R}^k$ with $\pi \circ s = Id_M$. The set of all sections is denoted by $\Gamma (M \times \mathbb{R}^k)$.    
\end{definition}

\begin{remark}
$\Gamma ( M \times \mathbb{R}^k) $ is naturally a $C^{\infty}(M)$-module. It can also be  naturally identified with the set of smooth functions $f : M \to \mathbb{R}^k$.  
\end{remark}

\begin{definition}

 Let $M$ be a smooth manifold and $k \in \mathbb{N}$. A \emph{vector bundle} of rank $k$ over $M$ is a smooth manifold $E$ equipped with a smooth map $\pi : E \to M$ and a collection of pairs $\{ ( U _i, \psi _i ) \} _{i \in I }$, called \emph{(local) trivializations}, such that:
 \begin{itemize}
     \item $\{ U_i \} _{i \in I }$ is an open cover of $M$. 
     \item For each $i$, $\psi _ i : \pi ^{-1} (U_i) \to U _i \times \mathbb{R}^k$ is a diffeomorphism with $\pi \circ \psi _i = \pi $, where we also denote by $\pi$ the projection onto the first coordinate $U_i \times \mathbb{R}^k \to U_i$.
     \item For each $x \in U_i \cap U_j$, the map 
     \[ \psi _j \circ \psi _i ^{-1} :  \{ x \} \times \mathbb{R}^k  \to \{ x \} \times \mathbb{R}^k  \] 
     is a linear isomorphism. 
 \end{itemize}
A \emph{section} of $E$ is a smooth map $s : M \to  E $ with $\pi \circ s = Id_M$. The set of all sections is denoted by $\Gamma (E)$.   
\end{definition}



\begin{remark}
    By the third property, for each $x \in M$, the preimage $E_x : = \pi ^{-1} (x)$ has a vector space structure inherited from the map $\psi_i $. Because of this, $\Gamma (E)$ is a $C^{\infty} (M)$-module. 
\end{remark}

\begin{example}
Let $M$ be an $n$-dimensional smooth manifold. Then the tangent bundle $\pi : TM \to M$ is a vector bundle of rank $n$.    
\end{example}


\begin{theorem}
        For any general construction in linear algebra, there is an analogous construction for vector bundles. This includes: Hom, the dual vector space, the direct sum, the tensor product, the tensor algebra, the exterior algebra, and the Clifford algebra.
\end{theorem}


\begin{example}
    For any vector bundles $E, F  $ over $M$,  there  are vector bundles  $E \oplus F ,$ $ E\otimes F,  $ $ Hom (E, F)$ over $M$ with
    \begin{gather*}
        (E \oplus F) _x =  E_x \oplus F_x  , \hspace{1.2cm} (E \otimes F) _x = E_x \otimes _{\mathbb{R}} F_x  , \\
          Hom(E,F)_x = Hom_{\mathbb{R}}(E_x, F_x)
    \end{gather*}
    for all $x \in M$.
\end{example}


\begin{solution}
We do the one of direct sum. The other two are similar, but the computations are a little more involved. Start by defining the set 
\[    E \oplus F : = \bigsqcup _{x \in M} (E _x \oplus F_x)        \]
and equip it with the map $\pi _{\oplus} : E \oplus F \to M$ that sends $(E _x \oplus F_x) $ to $x$. We now construct trivializations using the ones of $E$ and $F$. 

Take local trivializations $(U , \psi )$ and $(U , \zeta )$ of $E $ and $F$, respectively. That is, $U \subset M$ is an open set, and 
\begin{gather*}
        \psi : \pi_E^{-1} (U) \to U \times \mathbb{R}^k  \\
        \zeta : \pi_F^{-1} (U) \to U \times \mathbb{R}^{\ell} 
\end{gather*} 
diffeomorphisms, such that for each $x \in U$, the restrictions $\psi : E_x \to \{ x \} \times \mathbb{R}^k$ and $\zeta : F_ x \to \{ x \} \times \mathbb{R}^{\ell}$ are linear isomorphisms.  Then define 
\[   \psi \oplus \zeta : \pi _{\oplus}^{-1} (U) \to U \times \mathbb{R} ^k \times \mathbb{R}^{\ell}    \]
as
 \[  ( \psi \oplus \zeta )( v, w )  : = (x , \psi _2 (v), \zeta _2 (w) ) ,                     \]
where $v \in E_x$, $w \in F_x$, and $\psi _2$, $\zeta _2$ are the second coordinates of $\psi$ and $\zeta$, respectively. For each $x \in U$, the map $\psi \oplus \zeta : (E _x \oplus F_x) \to \{ x \} \times \mathbb{R}^{k + \ell }$ is a linear isomorphism. However, we need to check that when we consider other trivialization, the change of coordinates is smooth. 

Take $(V \overline{\psi}) $, $(V, \overline{\zeta})$ two local trivializiations of $E$ and $F$. We need to check that the change of coordinates 
\[  (\psi \oplus \zeta ) \circ (\overline{\psi} \oplus \overline{\zeta})^{-1} : (\overline{\psi} \oplus \overline{\zeta}) (\pi_{\oplus} ^{-1} (U \cap V) ) \to ( \psi \oplus \zeta ) (\pi_{\oplus} ^{-1} (U \cap V) )                       \]
is smooth. This happens to be very easy.  Note that 
\[  (\overline{\psi} \oplus \overline{\zeta}) (\pi_{\oplus} ^{-1} (U \cap V) ) =  ( \psi \oplus \zeta ) (\pi_{\oplus} ^{-1} (U \cap V) )   =   (U \cap V)  \times \mathbb{R}^k \times \mathbb{R}^{\ell}   .  \]
Then
\begin{align*}
    (\psi \oplus \zeta ) & \circ (\overline{\psi} \oplus \overline{\zeta})^{-1}  (  x, \lambda_1, \ldots , \lambda_k , \mu _1, \ldots , \mu _{\ell} )   \\
    & = (\psi \oplus \zeta ) ( \overline{\psi} ^{-1} (x, \lambda_1, \ldots , \lambda _k), \overline{\zeta}^{-1}(x, \mu_1, \ldots , \mu_{\ell}))    \\
    & = (  x , \psi _2 ( \overline{\psi}^{-1} (x, \lambda_1, \ldots , \lambda _k) ) , \zeta _2 (\overline{ \zeta }^{-1} ( x, \mu _1, \ldots , \mu _ {\ell}   ) ) )    .
\end{align*}
Since $\psi $ and $\overline{\psi}$ are diffeomorphisms, the second coordinate depends smoothly on $x$ and the $\lambda_i$'s. Since $\zeta $ and $\overline{\zeta}$ are diffeomorphisms, the third coordinate depends smoothly on $x$ and the $\mu _i $'s.     
\end{solution}




\begin{example}
    For any vector bundle $E \to M$,  there is a vector bundle  $E^{\ast} \to M$  with $ ( E^{\ast} )_x  = (E_x)^{\ast}$ for all $x \in M$. 
\end{example}


\begin{solution}
Start by defining the set
\[     E^{\ast} : = \bigsqcup_{x \in M} (E_x)^{\ast}      \]
and equip it with the map $\pi_{E^{\ast}}:  E ^{\ast} \to M$ that sends $(E_x)^{\ast}$ to $x$. We now construct trivializations using the ones of $E$.  


Let $(U, \psi) $ be a local trivialization of $E$. That is, $U \subset M$ is an open set and $\psi : \pi_E ^{-1} (U ) \to U \times \mathbb{R}^k$ a diffeomorphism where $\psi : E_x \to \{ x \} \times \mathbb{R}^k$ is a linear isomorphism for all $x \in U$. For each $x \in U$, define $v_i  (x) = \psi ^{-1} (x,e_i)$, where $e_i = (0, \ldots , 1, \ldots , 0)$ is the $i$-th canonical vector. Let $\{ \eta ^1 (x), \ldots , \eta ^k(x) \} \subset E_x^{\ast}$ be the dual basis. 





We can then define $\psi ^{\ast} : \pi _{E^{\ast}}^{-1}(U)  \to U \times \mathbb{R}^k$  by
    \[    \psi ^{\ast} \Big(  \sum a_i \eta ^i (x) \Big) = ( x , a_1, \ldots , a_k ).                 \]
For each $x \in U$, the map $\psi ^{\ast} : E_x^{\ast} \to \{ x \} \times  \mathbb{R}^k$  is a linear isomorphism. However, we need to check that when we consider other trivialization, the change of coordinates is smooth. 

Consider other local trivialization $(V, \zeta)$. We need to check that 
 \[\psi ^{\ast} \circ ( \zeta ^{\ast}  ) ^{-1} :  \zeta ^{\ast} (  \pi _{E^{\ast}}^{-1} (U \cap V) ) \to  \psi ^{\ast} (\pi _{E^{\ast}}^{-1} (U \cap V) ) )         \]
is smooth. Note that 
\[    \zeta ^{\ast} (  \pi _{E^{\ast}}^{-1} (U \cap V) ) = \psi ^{\ast} (\pi _{E^{\ast}}^{-1} (U \cap V) ) ) = (U \cap V ) \times \mathbb{R}^k  .        \]
Let $w_j (x) : = \zeta ^{-1} (x,e_j)$ and $\{ \alpha ^1 (x) , \ldots , \alpha ^k (x) \} \subset E_x ^{\ast}$ the dual basis. For each $x \in U \cap V$, find a matrix $M (x) = (m_{ij} (x))_{ij}$ such that
\[   w_j (x) = \sum _{i} m_{ij}(x) v_i (x)          .        \]
Since  $\psi $ and $\zeta$ are diffeomorphisms, the change of charts is smooth, and each $m_{ij} : U \to \mathbb{R}$ is smooth.  If we denote the inverse matrix by $M^{-1}(x) = (m^{ij}(x))_{ij}$, we have (ommiting the dependence on $x$):
\[     \sum _{i} m^{a i } \eta ^{i} (w_{b}) =  \sum _{i , \ell } m^{a i } \eta ^{i} (m_{\ell b} v_{\ell} ) = \sum_i  m^{ai} m_{ib} = \delta ^a_b ,                         \]
where $\delta ^a _b $ equals $1$ if $a = b$ and $0$ if $a \neq b$.  Therefore
\[     \alpha ^a = \sum _i m^{ai}\eta ^i .      \]
Consequently, 
\begin{align*}
    \psi ^{\ast} (  (\zeta ^{-\ast}&)  ^{-1}  (x, \lambda _1  , \ldots , \lambda _k )  )  = \psi ^{\ast} \left(   \sum _i \lambda _i \alpha ^i  (x)  \right) \\
    & = \psi ^{\ast} \left(   \sum _{i, \ell } \lambda _i  m^{i\ell } \eta  ^{\ell} (x)    \right)   \\
    & = \left( x, \sum_{i } \lambda_i m^{i 1} , \ldots , \sum_i \lambda _i m^{ik} \right). 
    \end{align*}
    Since the matrix $M^{-1}$ depends smoothly on $x$, the last expression depends smoothly on $x$ and the  $\lambda_i$'s. 
\end{solution}


\begin{example}
    The \emph{cotangent bundle} is the vector bundle $T^{\ast}M \to M$ with 
    \[ (T^{\ast}M)_x : =  T_x^{\ast}M    : = (T_xM)^{\ast}\text{  for all }x \in M. \]
\end{example}




\begin{definition}
    Let $E,F \to M$ be two vector bundles. A \emph{morphism} between $E$ and $F$ is a smooth map $\varphi : E \to F $ with $\pi \circ \varphi  = \pi $  such that $\varphi : E_x \to F_x$ is a linear map for each $x \in M$.     
    
    We say that $\varphi$ is a \emph{vector bundle isomorphism} if it is also a diffeomorphism, or equivalently, $\varphi : E_x \to F_x$ is a linear isomorphism for all $x \in M$. In such a case, we say that $E$ and $F$ are \emph{isomorphic} and we write $E \cong F$.
\end{definition}

\begin{theorem}
     For any natural isomorphism in linear algebra, there is an analogous isomorphism for vector bundles. 
\end{theorem}


\begin{example}
     For any vector bundle $E$ over $M$, we have
    \[    E  \cong E^{\ast \ast }    \]
\end{example}

\begin{example}
    For any vector bundles $E_1, E_2, E_3  $ over $M$, the direct sum and tensor product are commutative, associative, and distributive:
    \begin{gather*}
    E_1 \oplus E_2 \cong E_2 \oplus E_1 , \hspace{3cm} E_1 \otimes E_2 \cong E_2 \otimes E_1 , \\
       ( E_1 \oplus E_2) \oplus E_3 \cong E_1 \oplus (E_2 \oplus E_3) , \hspace{1cm} ( E_1 \otimes E_2) \otimes E_3 \cong E_1 \otimes (E_2 \otimes E_3) ,  \\
       ( E_1 \oplus E_2) \otimes E_3 \cong ( E_1 \otimes E_2 ) \oplus (E_2 \otimes E_3) .
    \end{gather*}
\end{example}



\begin{theorem}
For any vector bundles $E, F $ over $M$, there are natural isomorphisms
\begin{gather*}
     \Gamma (E \oplus F) \cong \Gamma (E) \oplus \Gamma (F)  ,  \\
      \Gamma (E \otimes F) \cong \Gamma (E) \otimes _{C^{\infty}(M)} \Gamma (F)  , \\
       \Gamma ( \, Hom ( E, F) \, ) \cong  Hom_{C^{\infty}(M)}( \, \Gamma (E) \, , \,  \Gamma (F) \, ) .
       \end{gather*}
In particular, we have  $    \Gamma (E^{\ast}) \cong    \Gamma (E) ^{\ast}  $, where the second dual is taken as a $C^{\infty}(M)$-module.
\end{theorem}

\begin{solution}
    See for example Conlon Lawrence. Differentiable manifolds, a first course. Birkh\"auser, 1993.
\end{solution}




















\end{document}